\begin{thebibliography}{99}

\bibitem{problem}
武汉理工大学数学建模训练赛命题组. 第9题：多区域算力协同与算电联合调度问题及附件数据说明, 2026.

\bibitem{wu_survey_2025}
Wu Y, Tang S, Yu C, Yang B, Sun C, Xiao J, Wu H, Feng J. Task Scheduling in Geo-Distributed Computing: A Survey. IEEE Transactions on Parallel and Distributed Systems, 2025, 36(10): 2073--2088. DOI: 10.1109/TPDS.2025.3591010.

\bibitem{guo_integrated_2025}
Guo H, Yu H, Wang M, Liu C, Li C. Integrated management of workloads and energy system for data centers. Energy, 2025, 327: 136400. DOI: 10.1016/j.energy.2025.136400.

\bibitem{taddy_marked_2012}
Taddy M A, Kottas A. Mixture Modeling for Marked Poisson Processes. Bayesian Analysis, 2012, 7(2): 335--362. DOI: 10.1214/12-BA711.

\bibitem{zou_poissonity_2022}
Zou Q, Zhu Y, Tan Y, Chen W. Diagnosing the coexistence of Poissonity and self-similarity in memory workloads. Journal of Network and Computer Applications, 2022, 205: 103455. DOI: 10.1016/j.jnca.2022.103455.

\bibitem{cohen_dac_2022}
Cohen M C, Zhang R, Jiao K. Data Aggregation and Demand Prediction. Operations Research, 2022, 70(5): 2597--2618. DOI: 10.1287/opre.2022.2301.

\bibitem{bertani_jbu_2025}
Bertani N, Jensen S T, Satopaa V A. Joint Bottom-up Method for Probabilistic Forecasting of Hierarchical Time Series. Operations Research, 2025, 73(6): 3260--3277. DOI: 10.1287/opre.2022.0113.

\bibitem{yang_repta_2026}
Yang L, Shahidehpour M, Chen X, Wang C, Fang X, Yang Q. High-speed REPTA algorithm for performance cost optimization in data centers considering workload delay tolerances. Applied Energy, 2026, 404: 127156. DOI: 10.1016/j.apenergy.2025.127156.

\bibitem{hanafy_carbonscaler_2023}
Hanafy W A, Liang Q, Bashir N, Irwin D, Shenoy P. CarbonScaler: Leveraging Cloud Workload Elasticity for Optimizing Carbon-Efficiency. Proceedings of the ACM on Measurement and Analysis of Computing Systems, 2023, 7(3): Article 57, 1--28. DOI: 10.1145/3626788.

\bibitem{zhang_dc_bess_2025}
Zhang Y, Tang H, Li H, Wang S. Unlocking the flexibilities of data centers for smart grid services: Optimal dispatch and design of energy storage systems under progressive loading. Energy, 2025, 316: 134511. DOI: 10.1016/j.energy.2025.134511.

\bibitem{vaicys_convex_bess_2024}
Vaičys J, Gudžius S, Jonaitis A, Račkienė R, Blinov A, Peftitsis D. A case study of optimising energy storage dispatch: Convex optimisation approach with degradation considerations. Journal of Energy Storage, 2024, 97(B): 112941. DOI: 10.1016/j.est.2024.112941.

\bibitem{wang_exact_relax_2024}
Wang Q, Wu W, Lin C, Xu S, Wang S, Tian J. An exact relaxation method for complementarity constraints of energy storages in power grid optimization problems. Applied Energy, 2024, 371: 123592. DOI: 10.1016/j.apenergy.2024.123592.

\bibitem{niu_benders_2021}
Niu T, Hu B, Xie K, Pan C, Jin H, Li C. Spacial coordination between data centers and power system considering uncertainties of both source and load sides. International Journal of Electrical Power \& Energy Systems, 2021, 124: 106358. DOI: 10.1016/j.ijepes.2020.106358.

\bibitem{ji_marginal_2021}
Ji K, Zhang F, Chi C, Song P, Zhou B, Marahatta A, Liu Z. A joint energy efficiency optimization scheme based on marginal cost and workload prediction in data centers. Sustainable Computing: Informatics and Systems, 2021, 32: 100596. DOI: 10.1016/j.suscom.2021.100596.

\bibitem{benders_1962}
Benders J F. Partitioning procedures for solving mixed-variables programming problems. Numerische Mathematik, 1962, 4: 238--252. DOI: 10.1007/BF01386316.

\bibitem{dantzig_wolfe_1960}
Dantzig G B, Wolfe P. Decomposition principle for linear programs. Operations Research, 1960, 8(1): 101--111. DOI: 10.1287/opre.8.1.101.

\end{thebibliography}
